\section{Permanente: dos escalados y un menor}
Según C, Lemas 7.2--7.3, $P_{ij}(1-P_{ij})$ es el producto de una función de fila
y otra de columna en el soporte de $A$. Escalando esta matriz se obtiene una matriz
bistocástica $D$ que conserva el soporte. Sus multiplicadores logarítmicos, en el gauge
de sumas iguales, forman $z_P=Wc+O_\infty(k^{-2})$: el defecto en el origen es $-c$,
el jacobiano difiere de $H$ en $O(k^{-1})$ en norma de sumas por fila,
y el Lema 7.3 da una solución $O_\infty(k^{-1})$. Por Taylor,
$D_{ij}=P_{ij}(1-P_{ij})\exp((z_P)_i+(z_P)_{k+j})$. Se tiene
$D-P=O_{\max}(k^{-2})$, sumas absolutas por línea $O(k^{-1})$ y brecha uniforme.
Cada permutación del soporte recibe el mismo peso en $D$. Por ello
$p_b=D_aD_b\per F/\per D$, donde $F$ es el menor de $D$.

Sólo en esta sección y en la siguiente, hasta la reconversión final, $U,R,\ell$
se forman con $D$. Denotemos por $S$ los cuatro extremos borrados, $T$ los
retenidos, $d=u_a+u_b$, $m=k-2$, y por $h$ la masa de $D$ en el rectángulo borrado
de dos filas por dos columnas. Los defectos del menor, su jacobiano y la corrección
de primer orden son
\[
 g=(\mathbf1-F\mathbf1,\mathbf1-F^T\mathbf1),\quad
 H_F=\begin{pmatrix}\diag(F\mathbf1)&F\\F^T&\diag(F^T\mathbf1)\end{pmatrix},
 \qquad z_0=H_F^+g .
\]
En el complemento de $(\mathbf1,-\mathbf1)$, $H_F$ tiene brecha uniforme.
En efecto $F=G+\mathbf1\mathbf1^T/k$ donde $G$ es la submatriz de $D-J_k$;
$\|G\|\le r'<1$, las diagonales difieren de uno en $O(1/k)$ y la contribución
de la matriz constante a la forma en ese gauge es no negativa. Además $H_F^+$
está acotada en norma infinito: si $H_F z=b$ en el gauge, entonces
$\|z\|_2\le C\sqrt{k}\|b\|_\infty$; cada fila o columna de $F$ tiene norma
euclídea $O(k^{-1/2})$, y las dos ecuaciones por bloques dan la cota afirmada.
El operador extendido proyecta sobre el gauge con norma infinito acotada.

Para escalar $F$, el mapa de márgenes logarítmicos tiene jacobiano $H_F$ en cero;
en una bola de radio $r_0$ su variación en norma infinito es $O(r_0)$ por la
suma de cada línea. El punto fijo con inversa congelada existe por contracción
en la bola de radio $2\|H_F^+g\|_\infty=O(1/k)$ (si $g=0$, tomar cero).
Su vector $z=(u,v)$ satisface $z=z_0+O_\infty(k^{-2})$. La matriz
$D'=e^{u_i+v_j}F_{ij}$ es bistocástica con el mismo soporte,
$\max D'_{ij}=O(1/k)$ y $\|D'-J_m\|\le r'+O(1/k)<1$:
$F-J_m=G+O_{\rm op}(1/k)$ y el escalado cambia $F$ en $O_{\rm op}(1/k)$.
Así el menor mismo admite la aproximación del permanente. No se presupone que
su punto estacionario de torres sea la restricción del del padre.

En $T$ se tiene $g=2\mathbf1/k+(Ud)_T$ y
\[
 z_0=\mathbf1/k-(Rd)_T+O_\infty(k^{-2}).
\]
Para verificar la segunda igualdad, la parte de $H_F$ igual a $(I+U)_{TT}$
actúa sobre el candidato como $\mathbf1/k+(Ud)_T+O_\infty(k^{-2})$ por
$(I+U)R=-U$. La parte constante cruzada añade $\mathbf1/k+O_\infty(k^{-2})$:
las sumas completas de $Rd$ en cada bloque son cero y sólo se quitaron cuatro
entradas $O(1/k)$. La perturbación diagonal $-\diag g$ cuesta $O_\infty(k^{-2})$.
Ajustar el gauge cambia cada coordenada en $O(k^{-2})$. Extendiendo las sumas
con error $O(k^{-2})$ en los productos escalares se obtiene
\[
 g^Tz_0=4/k-d^TURd+O(k^{-2})=2h-4/k+d^TRd+O(k^{-2}).
\]
Aquí $UR=-U-R$, $d^TUd=2h-8/k$ y $U,R$ anulan los modos constantes.
La identidad de masa $m=\sum F_{ij}e^{u_i+v_j}$, desarrollada hasta grado dos,
da $\sum z=2-h+g^Tz-\tfrac12 z^T H_F z+O(k^{-2})$:
$\sum F=k-4+h$ y el resto cúbico cuesta $O(k\|z\|_\infty^3)$.
Por tanto $\sum z=2-h+\tfrac12g^Tz_0+O(k^{-2})$. Es el logaritmo del factor
por el que se multiplica $\per F$ al escalarlo.
